% Modelo de TCC em LaTeX - PROFMAT
\documentclass[12pt,a4paper]{report}
\usepackage[utf8]{inputenc}
\usepackage[brazil]{babel}
\usepackage{amsmath, amssymb, amsfonts}
\usepackage{graphicx}
\usepackage{geometry}
\usepackage{setspace}
\usepackage{listings}
\usepackage{color}
\usepackage{hyperref}

\geometry{a4paper, margin=2.5cm}
\onehalfspacing

\definecolor{lightgray}{gray}{0.95}
\lstset{
	backgroundcolor=\color{lightgray},
	basicstyle=\ttfamily\footnotesize,
	breaklines=true,
	frame=single,
	numbers=left,
	numberstyle=\tiny, 
	stepnumber=1,
	tabsize=4,
	language=Python,
	showstringspaces=false
}

\begin{document}
	
	% CAPA
	\begin{center}
		\Large{\textbf{UNIVERSIDADE XXXXXXX}} \\
		\large{PROGRAMA DE MESTRADO PROFISSIONAL EM MATEMÁTICA - PROFMAT} \\
		\vspace{3cm}
		\LARGE{\textbf{O uso do Python no ensino de métodos numéricos discretos no Ensino Médio}} \\
		\vspace{3cm}
		\large{Autor: Seu Nome Aqui} \\
		\large{Orientador: Prof. Dr. Nome do Orientador} \\
		\vfill
		\large{Cidade - Ano}
	\end{center}
	\newpage
	
	% RESUMO
	\begin{abstract}
		Este trabalho apresenta uma proposta de ensino de Matemática na Educação Básica utilizando a linguagem Python como ferramenta de apoio à exploração de métodos numéricos discretos. O projeto visa desenvolver o pensamento computacional e a compreensão de conceitos como aproximação, erro e convergência, por meio da resolução de equações e visualização gráfica de funções. A proposta está alinhada à BNCC, especialmente às habilidades EM13MAT401, EM13MAT402, EM13MAT501, EM13MAT502, EM13MAT503 e EM13MAT510, que incentivam o uso de tecnologias digitais, a análise de funções e a conversão entre representações algébricas e geométricas. \\ 
		\textbf{Palavras-chave:} Python; Métodos numéricos; Ensino Médio; BNCC; Equações; Aproximações.
	\end{abstract}
	\newpage
	
	% INTRODUÇÃO
	\chapter{Introdução}
	
	A integração entre Matemática e Computação representa uma oportunidade para tornar o ensino mais investigativo e dinâmico. A Base Nacional Comum Curricular (BNCC) destaca a importância do uso de tecnologias digitais e do pensamento computacional no desenvolvimento das competências matemáticas. Este trabalho propõe uma sequência didática utilizando Python para explorar métodos numéricos discretos aplicados à resolução de equações e funções, aproximando o estudante da experimentação científica e do raciocínio lógico.
	
	O problema central desta pesquisa é: \textit{como o uso de métodos numéricos discretos, implementados em Python, pode favorecer a compreensão de conceitos de funções, raízes e aproximações no Ensino Médio?}
	
	O objetivo geral é investigar o uso do Python como ferramenta de apoio ao ensino de Matemática, com foco nos métodos da Bisseção, Secante e Newton-Raphson. Especificamente, pretende-se:
	\begin{itemize}
		\item Implementar algoritmos simples de métodos numéricos discretos;
		\item Visualizar graficamente o comportamento das funções e suas raízes;
		\item Analisar o erro e a convergência dos métodos iterativos;
		\item Relacionar os resultados obtidos à BNCC e às práticas pedagógicas investigativas.
	\end{itemize}
	
	\newpage
	\chapter{Fundamentação Teórica}
	
	\section{A BNCC e o ensino de Matemática}
	A BNCC (2018) propõe que o estudante desenvolva competências relacionadas à resolução de problemas, à modelagem e ao uso de tecnologias digitais. As habilidades EM13MAT401 a EM13MAT510 orientam o trabalho com funções, representações gráficas e linguagens de programação.
	
	\section{Métodos numéricos discretos}
	Os métodos numéricos buscam soluções aproximadas para equações que não podem ser resolvidas analiticamente. Entre os mais acessíveis para o Ensino Médio estão:
	
	\subsection{Método da Bisseção}
	\begin{equation}
		m = \frac{a + b}{2}
	\end{equation}
	O intervalo $[a,b]$ é atualizado conforme o sinal de $f(a)f(m)$, até que $|b-a| < \varepsilon$ ou $|f(m)| < \varepsilon$.
	
	\subsection{Método da Secante}
	\begin{equation}
		x_{n+1} = x_n - f(x_n) \cdot \frac{x_n - x_{n-1}}{f(x_n) - f(x_{n-1})}
	\end{equation}
	Este método dispensa o cálculo da derivada, utilizando apenas valores anteriores da função.
	
	\subsection{Método de Newton-Raphson}
	\begin{equation}
		x_{n+1} = x_n - \frac{f(x_n)}{f'(x_n)}
	\end{equation}
	Critérios de parada:
	\begin{equation}
		|x_{n+1} - x_n| < \varepsilon \quad \text{ou} \quad |f(x_{n+1})| < \varepsilon
	\end{equation}
	
	\subsection{Erro absoluto e relativo}
	\begin{align}
		E_a &= |x_{aprox} - x_{real}| \\
		E_r &= \frac{|x_{aprox} - x_{real}|}{|x_{real}|}
	\end{align}
	
	\newpage
	\chapter{Metodologia}
	
	A pesquisa adota uma abordagem qualitativa e exploratória, com elementos de pesquisa-ação. As atividades foram planejadas para turmas do Ensino Médio, utilizando o Google Colab como ambiente de programação.
	
	\section{Etapas da aplicação}
	\begin{enumerate}
		\item Introdução ao Python e às funções matemáticas básicas.
		\item Plotagem de gráficos de funções e identificação de raízes.
		\item Implementação dos métodos numéricos discretos.
		\item Análise de erro e convergência.
		\item Desenvolvimento de projeto final pelos alunos.
	\end{enumerate}
	
	\section{Exemplo de código Python}
	\begin{lstlisting}[language=Python, caption=Implementação do método da Bisseção]
		import math
		
		def bissecao(f, a, b, tol=1e-6, max_iter=100):
		fa, fb = f(a), f(b)
		if fa * fb > 0:
		raise ValueError('Intervalo inválido')
		for k in range(max_iter):
		m = (a + b) / 2
		fm = f(m)
		if abs(fm) < tol or (b - a)/2 < tol:
		return m, k
		if fa * fm < 0:
		b, fb = m, fm
		else:
		a, fa = m, fm
		return m, max_iter
		
		f = lambda x: x**2 - 3
		raiz, iteracoes = bissecao(f, 1, 2)
		print('Raiz aproximada:', raiz)
	\end{lstlisting}
	
	\newpage
	\chapter{Resultados Esperados}
	
	Espera-se que os alunos compreendam o conceito de aproximação, erro e convergência, além de desenvolverem raciocínio lógico e pensamento computacional. A metodologia visa fortalecer a relação entre teoria e prática e despertar o interesse pela modelagem matemática.
	
	\chapter{Considerações Finais}
	
	O ensino de métodos numéricos discretos com Python permite trabalhar conteúdos de funções e equações de forma aplicada e investigativa, desenvolvendo habilidades previstas na BNCC e promovendo a integração entre Matemática e Computação.
	
	\chapter*{Referências}
	\begin{itemize}
		\item BRASIL. Base Nacional Comum Curricular. MEC, 2018.
		\item BORBA, M. C.; PENTEADO, M. G. Informática e Educação Matemática. Autêntica, 2016.
		\item COSTA, R. S. Ensino de Métodos Numéricos com Python: uma abordagem didática. Revista de Educação Matemática, 2020.
		\item VALENTE, J. A. O Computador na Sociedade do Conhecimento. Campinas: UNICAMP, 2018.
		\item WING, J. M. Computational Thinking. Communications of the ACM, 2006.
	\end{itemize}
	
\end{document}